\documentclass[11pt,a4paper]{article}
\usepackage[utf8]{inputenc}
\usepackage[T1]{fontenc}
\usepackage[brazil]{babel}
\usepackage[margin=2.4cm]{geometry}
\usepackage{booktabs}
\usepackage{amsmath}
\usepackage{amssymb}
\usepackage{xcolor}
\usepackage{hyperref}
\hypersetup{colorlinks=true, linkcolor=blue!50!black, urlcolor=blue!50!black}

\newcommand{\keff}{k_{\text{eff}}}
\newcommand{\aprovado}{\textcolor{green!45!black}{\textbf{APROVADO}}}
\newcommand{\reprovado}{\textcolor{red!70!black}{\textbf{REPROVADO}}}
\newcommand{\e}[1]{\ensuremath{\times 10^{#1}}}

\title{Relatório de Avaliação Independente\\
\large Solver de Difusão de Nêutrons Monoenergética 1D/2D\\
(Opções 1 e 3 do projeto final de Física de Reatores)}
\author{Bateria de testes de verificação, física, consistência numérica e robustez}
\date{6 de julho de 2026}

\begin{document}
\maketitle

\begin{abstract}
Este relatório documenta uma bateria independente de 25 verificações aplicadas ao
solver \texttt{difusao} (volumes finitos \emph{cell-centered}, LU esparsa/CG,
\emph{power iteration}/Wielandt), organizada em quatro frentes: (A)~verificação
contra soluções analíticas, (B)~física de reatores, (C)~consistência numérica e
(D)~robustez de entrada. Resultado global: \textbf{a física e a numérica do
código estão sólidas} --- 18/18 testes da suíte própria passam, os 10 casos
obrigatórios conservam nêutrons com resíduo $\leq 3{,}5\e{-13}$, o autovalor
converge para o benchmark analítico com erro $8{,}5\e{-12}$ após extrapolação de
Richardson, e todas as invariâncias físicas são satisfeitas em precisão de
máquina. Foram identificadas \textbf{5 falhas de robustez na camada de entrada}
(tratamento de erros e avisos ao usuário), nenhuma delas afetando a correção dos
resultados numéricos. As cinco foram \textbf{corrigidas} com alterações
localizadas em \texttt{input\_data.py} e \texttt{malha.py}
(Seção~\ref{sec:correcoes}); após as correções, a bateria completa e a suíte
própria (18/18) passam integralmente. Todos os números deste relatório são
reprodutíveis via \texttt{relatorio/testes\_professor.py}.
\end{abstract}

\section{Objeto avaliado e metodologia}

O código resolve a equação de difusão monoenergética estacionária
\begin{equation}
-\nabla\!\cdot\!\bigl(D\,\nabla\phi\bigr) + \Sigma_a\,\phi
  = Q \quad\text{ou}\quad \frac{1}{\keff}\,\nu\Sigma_f\,\phi,
\qquad D = \frac{1}{3\Sigma_t},
\end{equation}
em meios heterogêneos 1D/2D, por volumes finitos \emph{cell-centered} com média
harmônica de $D$ nas interfaces, condição de vácuo de Marshak
($\alpha = 2D/(h+4D)$) e reflexão por corrente nula. Sistemas lineares por LU
esparsa (\texttt{splu}, MMD) ou CG+Jacobi; autovalor por \emph{power iteration}
ou Wielandt em dois estágios.

A avaliação exercita o código \emph{sem alterá-lo}: os módulos são importados
diretamente (testes A--C) ou o executável \texttt{main.py} é invocado como o
usuário final o faria (testes D). Critérios de aprovação são declarados
\emph{a priori} em cada teste.

\subsection*{Ambiente}
Python 3.14.5 (venv), \texttt{numpy}/\texttt{scipy}/\texttt{pyyaml} conforme
\texttt{requirements.txt}, Linux x86-64.
\textbf{Observação de reprodutibilidade:} o interpretador do sistema não possuía
as dependências instaladas; foi necessário criar um ambiente virtual. O
\texttt{README} instrui corretamente o \texttt{pip install}, mas recomenda-se
declarar explicitamente a criação do venv.

\section{Suíte de testes própria e casos obrigatórios}

A suíte \texttt{pytest} do próprio projeto foi executada integralmente:
\textbf{18/18 testes passam} em $0{,}3$\,s. Em seguida os 10 casos obrigatórios
do enunciado (placa de 20\,cm, $\Sigma_t = 1$, vácuo em ambos os lados):

\begin{table}[h!]\centering\small
\begin{tabular}{lccccc}
\toprule
Caso & $\Sigma_{s1}$ & $\Sigma_{s2}$ & Absorção & Fuga & Resíduo do balanço \\
\midrule
01 & 0{,}8    & 0{,}8   & 9{,}148749 & 0{,}851251 & $2{,}03\e{-14}$ \\
02 & 0{,}8    & 0{,}9\textsuperscript{*}  & 9{,}072901 & 0{,}927099 & $2{,}04\e{-14}$ \\
03 & 0{,}8    & 0{,}99  & 8{,}668636 & 1{,}331364 & $1{,}03\e{-14}$ \\
04 & 0{,}8    & 0{,}999\textsuperscript{*} & 8{,}525411 & 1{,}474589 & $7{,}64\e{-15}$ \\
05 & 0{,}9    & --      & 6{,}155659 & 3{,}844341 & $1{,}41\e{-13}$ \\
06 & 0{,}99   & --      & 5{,}107640 & 4{,}892360 & $1{,}60\e{-13}$ \\
07 & 0{,}99   & --      & 4{,}930043 & 5{,}069957 & $1{,}41\e{-13}$ \\
08 & 0{,}999  & --      & 2{,}835959 & 7{,}164041 & $1{,}92\e{-13}$ \\
09 & 0{,}9999 & --      & 4{,}756272 & 5{,}243728 & $2{,}36\e{-13}$ \\
10 & 0{,}9999 & --      & 2{,}464999 & 7{,}535001 & $3{,}50\e{-13}$ \\
\bottomrule
\end{tabular}
\caption{Balanço de nêutrons dos 10 casos obrigatórios (fonte total $=10$ em
todos). Combinações conforme os arquivos \texttt{inputs/caso01--10.yaml}
(colunas $\Sigma_s$ meramente indicativas; ver os arquivos para a definição
completa). Todos os resíduos ficam 5--6 ordens de grandeza abaixo do limiar de
aceitação do próprio código ($10^{-8}$). \aprovado}
\end{table}

Os exemplos adicionais também executam corretamente: 2D com LU
(resíduo $7{,}4\e{-14}$), 2D com CG+Jacobi (297 iterações, resíduo
$1{,}8\e{-12}$) e autovalor Wielandt ($\keff = 1{,}27764735$, 34 iterações,
3~fatorações LU, resíduo $2{,}9\e{-16}$).

\section{Frente A --- Verificação contra soluções analíticas}

\subsection{A1. Meio infinito homogêneo}
Reflexão em ambos os lados com fonte uniforme deve reproduzir
$\phi = Q/\Sigma_a = 2{,}0$ exatamente. Erro relativo máximo obtido:
$2{,}5\e{-14}$ (precisão de máquina). \hfill \aprovado

\subsection{A2. Ordem de convergência espacial (fonte fixa)}
Meia-placa (reflexiva--vácuo) contra a solução analítica
$\phi(x) = \tfrac{Q}{\Sigma_a}\bigl[1 - \cosh(x/L_d)/\text{den}\bigr]$:

\begin{center}\small
\begin{tabular}{lccccc}
\toprule
$N$ & 40 & 80 & 160 & 320 & 640 \\
\midrule
erro $\max$ & $1{,}56\e{-3}$ & $4{,}17\e{-4}$ & $1{,}08\e{-4}$ & $2{,}73\e{-5}$ & $6{,}88\e{-6}$ \\
ordem local & --- & 1{,}907 & 1{,}955 & 1{,}978 & 1{,}989 \\
\bottomrule
\end{tabular}
\end{center}
A ordem tende monotonicamente a 2, como esperado para o esquema. \hfill \aprovado

\subsection{A3. Placa nua crítica (benchmark analítico de autovalor)}
\textbf{Este teste não existia na suíte do projeto} --- a suíte original só
verificava Wielandt $\equiv$ \emph{power} (consistência interna, não
verificação). Para a placa nua ($L=20$\,cm, $\Sigma_t=1$, $\Sigma_s=0{,}9$,
$\nu\Sigma_f=0{,}15$) com condição de Marshak, o modo fundamental satisfaz
\begin{equation}
\tan\!\Bigl(\frac{BL}{2}\Bigr) = \frac{1}{2DB},
\qquad
\keff^{\text{ana}} = \frac{\nu\Sigma_f}{\Sigma_a + DB^2},
\end{equation}
cuja raiz dá $B = 0{,}14729157$\,cm$^{-1}$ e
$\keff^{\text{ana}} = 1{,}3988413535$.

\begin{center}\small
\begin{tabular}{lccccc}
\toprule
$N$ & 50 & 100 & 200 & 400 & 800 \\
\midrule
$|\keff - \keff^{\text{ana}}|$ & $3{,}07\e{-5}$ & $7{,}67\e{-6}$ & $1{,}92\e{-6}$ & $4{,}79\e{-7}$ & $1{,}20\e{-7}$ \\
ordem local & --- & 2{,}000 & 2{,}000 & 2{,}000 & 2{,}000 \\
\bottomrule
\end{tabular}
\end{center}
Convergência de ordem 2 impecável. A extrapolação de Richardson das duas malhas
mais finas dá $\keff = 1{,}3988413535$, com erro de $8{,}5\e{-12}$ contra o
valor analítico. O acoplamento fissão--fuga--absorção está correto.
\hfill \aprovado

\section{Frente B --- Física de reatores}

\subsection{B1. Continuidade da corrente na interface de material}
No caso $\Sigma_{s1}=0{,}99 \to \Sigma_{s2}=0{,}8$ (interface em $x=10$\,cm),
a corrente reconstruída pelos dois lados da interface:
$J^- = J^+ = 3{,}39206395$, salto relativo $1{,}4\e{-14}$. A média harmônica
garante $J^-=J^+$ por construção; a verificação numérica confirma a
implementação. \hfill \aprovado

\subsection{B2. Sensibilidade de $\keff$ (teoria de perturbação)}
Em um grupo, multiplicar $\nu\Sigma_f$ por $1{,}01$ deve escalar $\keff$
\emph{exatamente} por $1{,}01$ (mesmo fluxo, operador de fissão escalado).
Obtido: razão $= 1{,}0100000000$, desvio $8{,}4\e{-14}$. Endurecer o refletor
($\Sigma_s\!: 0{,}95 \to 0{,}90$, i.e.\ $\Sigma_a$ maior) reduz $\keff$ de
$1{,}277574$ para $1{,}274863$ --- sinal físico correto. \hfill \aprovado

\subsection{B3. Efeito do refletor}
Núcleo nu de 20\,cm ($k = 1{,}269401$) versus o mesmo núcleo com 20\,cm de
refletor de cada lado ($k = 1{,}277574$): ganho de reatividade de
$+817$\,pcm e achatamento do fluxo no núcleo (fator de pico
$1{,}480 \to 1{,}321$). Ambos os efeitos na direção fisicamente correta.
\hfill \aprovado

\subsection{B4. Razão de dominância e Wielandt}
Dois núcleos físseis fracamente acoplados (razão de dominância
$\approx 0{,}9997$): \emph{power} converge em 128 iterações, Wielandt em 34
(3~fatorações LU), ganho de $3{,}8\times$, com
$|k_W - k_P| = 3{,}0\e{-9}$. O Wielandt entrega o mesmo modo fundamental com
custo muito menor exatamente no regime em que a \emph{power iteration} sofre.
\hfill \aprovado

\subsection{B5. Quantificação do limite difusivo (quantitativo)}
O código usa distância extrapolada de Marshak $d = 2D = 0{,}6667$\,cm; a
teoria de transporte dá $d = 0{,}7104\,\lambda_{tr} = 0{,}7104$\,cm. Comparando
as soluções analíticas de meia-placa com cada $d$:

\begin{center}\small
\begin{tabular}{lcc}
\toprule
 & $c = 0{,}8$ & $c = 0{,}9999$ \\
\midrule
diferença relativa máxima em $\phi$ & $1{,}44\,\%$ & $0{,}77\,\%$ \\
diferença relativa em $\phi$ na borda & $4{,}06\,\%$ & $6{,}14\,\%$ \\
\bottomrule
\end{tabular}
\end{center}
Ou seja: a limitação declarada no \texttt{README} é real, mas
\emph{quantificada}: o erro de modelo concentra-se na vizinhança da borda de
vácuo ($\sim$4--6\,\% no fluxo de borda) e é $\lesssim 1{,}5\,\%$ no interior.
É limitação da aproximação de difusão, não defeito de implementação.

\section{Frente C --- Consistência numérica}

\subsection{C1. Pior caso combinado: 2D + malha graduada + autovalor}
Configuração que nenhum teste da suíte original cobria simultaneamente:
domínio 2D $40\times40$\,cm, graduação $p_x=1{,}8$, $p_y=1{,}4$, núcleo
central físsil, BC mista (3 vácuos + 1 reflexiva). Resultado:
$\keff = 1{,}266363$, \textbf{resíduo de balanço $1{,}1\e{-16}$} em 14
iterações. A conservação sobrevive à combinação mais agressiva de
funcionalidades. \hfill \aprovado

\subsection{C2. Independência de malha e extrapolação de Richardson}
$\keff$ do reator refletido em malhas encaixadas:

\begin{center}\small
\begin{tabular}{lcccc}
\toprule
$N$ & 150 & 300 & 600 & 1200 \\
\midrule
$\keff$ & 1{,}27760317 & 1{,}27757966 & 1{,}27757377 & 1{,}27757229 \\
\bottomrule
\end{tabular}
\end{center}
Ordem observada em $k$: $1{,}999$. Extrapolado de Richardson:
$\keff^\infty = 1{,}27757180$; a malha de $N=1200$ difere disso por apenas
$4{,}9\e{-7}$ ($<0{,}05$\,pcm). \hfill \aprovado

\subsection{C3. Invariâncias geométricas}
Espelhar o problema 1D (trocar esquerda$\leftrightarrow$direita): diferença
relativa $4{,}3\e{-14}$. Transpor o problema 2D (rotação de $90^\circ$):
diferença $3{,}1\e{-15}$. Nenhuma assimetria espúria de discretização ou de
indexação. \hfill \aprovado

\subsection{C4. Sensibilidade à semente do chute inicial}
O autovalor usa chute aleatório com semente fixa (42). Com sementes
$\{1, 7, 42, 123, 999\}$, a dispersão de $\keff$ é $3{,}8\e{-14}$ ---
ordens de grandeza abaixo de \texttt{tol\_k} $=10^{-8}$. O critério duplo de
parada (em $k$ e na forma do fluxo) é efetivamente robusto ao chute.
\hfill \aprovado

\subsection{C5. Subcrítico profundo ($k \approx 0{,}39$)}
Preocupação levantada \emph{a priori}: os \emph{shifts} fixos do Wielandt
($\delta = 0{,}02/0{,}002$) assumem $k$ bem estimado nas 8 iterações livres.
Teste com $k_\infty = 0{,}4$: Wielandt converge em 15 iterações
($k = 0{,}39058487$), \emph{power} em 81, diferença $6{,}2\e{-9}$. A
estratégia de dois estágios se mostrou robusta também longe de $k=1$.
\hfill \aprovado

\subsection{C6. Malha mínima aceita ($N_x = 3$)}
Mesmo na malha mais grosseira que a validação permite, o balanço fecha com
resíduo $1{,}8\e{-16}$ (a \emph{acurácia} é obviamente baixa, mas a
\emph{conservação} --- propriedade estrutural do esquema FV --- é exata).
\hfill \aprovado

\section{Frente D --- Robustez da camada de entrada}

Testes executados contra \texttt{main.py} como usuário final. Critério:
entrada inválida deve produzir \textbf{erro limpo} (mensagem
\texttt{[ERRO de input]}, sem \emph{traceback}) ou, quando a entrada é
tecnicamente válida porém suspeita, um \textbf{aviso}.

\begin{table}[h!]\centering\small
\begin{tabular}{p{4.6cm}p{5.6cm}l}
\toprule
Teste & Comportamento observado & Veredito \\
\midrule
D1: $\Sigma_s > \Sigma_t$ & Erro limpo com mensagem clara e valores ecoados & \aprovado \\
\addlinespace
D2: regiões não cobrem o domínio & \texttt{ValueError} com \emph{traceback} cru (a mensagem em si é boa, mas vaza pilha: \texttt{main.py} só captura \texttt{ErroDeInput}) & \reprovado \\
\addlinespace
D3: JSON malformado & \texttt{JSONDecodeError} com \emph{traceback} cru & \reprovado \\
\addlinespace
D4: \texttt{Nx: "quatrocentos"} & \texttt{ValueError} de \texttt{int()} com \emph{traceback} cru, sem apontar o campo & \reprovado \\
\addlinespace
D5: \texttt{fonte} $>0$ em modo autovalor & Executa e ignora a fonte \textbf{silenciosamente} (código de saída 0, nenhum aviso) & \reprovado \\
\addlinespace
D6: arquivo YAML vazio & Erro limpo (``o input deve ser um mapeamento'') & \aprovado \\
\addlinespace
D7: interface de material coincidindo com centróide de célula & O resultado depende da \textbf{ordem} em que as regiões são listadas: diferença relativa de $68{,}7\,\%$ entre as duas ordens, sem qualquer aviso & \reprovado \\
\bottomrule
\end{tabular}
\caption{Robustez de entrada \emph{antes das correções}: 3 aprovados (D1, D6 e
o já coberto pela suíte), 5 reprovados. Nenhuma falha compromete resultados de
entradas válidas. Ver Seção~\ref{sec:correcoes} para o estado final.}
\end{table}

\paragraph{Análise das falhas.} Todas as cinco são da mesma família ---
\emph{tratamento de borda da camada de entrada} --- e têm correção localizada:
\begin{enumerate}
  \item \textbf{D2/D4}: o \texttt{ValueError} de \texttt{malha.py} e o de
        \texttt{int()} não são \texttt{ErroDeInput}, então escapam do
        \texttt{except} de \texttt{main.py:106}. Converter/encapsular esses
        erros resolve D2 e D4 de uma vez.
  \item \textbf{D3}: capturar \texttt{json.JSONDecodeError} e
        \texttt{yaml.YAMLError} em \texttt{ler\_input} e reemitir como
        \texttt{ErroDeInput} com linha/coluna.
  \item \textbf{D5}: um aviso de uma linha no eco do input
        (``fonte fixa ignorada em modo autovalor'') basta.
  \item \textbf{D7}: é o caso mais sutil e o mais perigoso, porque é
        \emph{silencioso e fisicamente grande} (68\,\% no caso construído,
        deliberadamente patológico: $N=10$ com interface exatamente sobre um
        centróide). Mitigações possíveis: detectar interface não alinhada com
        faces da malha e avisar, ou fazer \emph{snap} das interfaces às faces.
        O \texttt{README} documenta a regra ``a última região sobrepõe'', mas
        documentação não substitui aviso em tempo de execução.
\end{enumerate}

\section{Correções aplicadas}
\label{sec:correcoes}

As cinco falhas da frente D foram corrigidas, sem alterar uma linha do núcleo
numérico (\texttt{montagem.py}, \texttt{solver.py}, \texttt{balanco.py}
intactos):

\begin{enumerate}
  \item \textbf{D2} --- \texttt{malha.py} passa a lançar \texttt{ErroDeInput}
        (subclasse de \texttt{ValueError}, preservando a compatibilidade da
        suíte de testes) quando as regiões não cobrem o domínio; o erro agora é
        capturado limpo por \texttt{main.py}.
  \item \textbf{D3} --- \texttt{ler\_input} captura
        \texttt{json.JSONDecodeError} e \texttt{yaml.YAMLError} e reemite como
        \texttt{ErroDeInput} com a posição do defeito no arquivo.
  \item \textbf{D4} --- novos validadores \texttt{\_numero}/\texttt{\_inteiro}
        exigem tipo numérico em \texttt{Lx}, \texttt{Nx}, \texttt{Ly},
        \texttt{Ny}, apontando o campo ofensor na mensagem.
  \item \textbf{D5} --- em modo autovalor com \texttt{fonte} $> 0$, é emitido
        \texttt{[AVISO]} explícito de que a fonte fixa é ignorada.
  \item \textbf{D7} --- \texttt{materiais\_1d}/\texttt{materiais\_2d} detectam
        borda de região coincidente com centróide de célula e avisam que a
        atribuição segue a última região listada, recomendando alinhar a
        interface a uma face da malha.
\end{enumerate}

Após as correções, a bateria completa foi reexecutada:
\textbf{21/21 verificações aprovadas} (frentes A--D), suíte própria 18/18, e os
12 casos de \texttt{inputs/} executam sem regressão (mesmos resíduos de
balanço). Na mesma revisão foram removidos comentários informais do
\texttt{README} e comentários supérfluos do código e dos arquivos de entrada,
mantendo apenas os \emph{docstrings}-título de cada função.

\section{Veredito}

\begin{center}
\begin{tabular}{lccl}
\toprule
Frente & Testes & Aprovados & Situação \\
\midrule
Suíte própria + casos obrigatórios & 18 + 13 & todos & sólida \\
A. Verificação analítica & 3 & 3/3 & sólida \\
B. Física de reatores & 5 & 5/5 & sólida \\
C. Consistência numérica & 6 & 6/6 & sólida \\
D. Robustez de entrada (antes) & 7 & 2/7 & \textbf{frágil} \\
D. Robustez de entrada (após correções) & 7 & 7/7 & sólida \\
\bottomrule
\end{tabular}
\end{center}

\textbf{Núcleo numérico e física: aprovados sem ressalvas.} O código
demonstra ordem 2 verificada contra dois benchmarks analíticos independentes
(fonte fixa \emph{e} autovalor --- este último ausente da suíte original e
adicionado por esta avaliação), conservação de nêutrons em precisão de máquina
mesmo no pior caso combinado, sensibilidades com o sinal e a magnitude
corretos, invariâncias exatas e solvers consistentes entre si inclusive em
regime subcrítico profundo.

\textbf{Camada de entrada: falhas identificadas e corrigidas.} As cinco falhas
D2--D5 e D7 eram todas de \emph{experiência de erro}, não de física; nenhuma
produzia resultado errado a partir de entrada válida e bem posta. Após as
correções da Seção~\ref{sec:correcoes}, os sete testes da frente D passam.

\textbf{Pendências documentais:} o \texttt{README} referencia notebooks
(\texttt{difusao\_opcoes\_1\_e\_3.ipynb}) ausentes do repositório e contém
comentários informais que não condizem com a qualidade do restante do
trabalho; as dependências exigem venv não documentado explicitamente.

\appendix
\section{Reprodutibilidade}
\begin{itemize}
  \item Suíte própria: \texttt{.venv/bin/python -m pytest tests/ -v}
        (saída em \texttt{relatorio/dados/pytest\_saida.txt});
  \item Casos obrigatórios: \texttt{.venv/bin/python main.py inputs/caso*.yaml}
        (\texttt{relatorio/dados/casos\_obrigatorios.txt});
  \item Bateria A--D: \texttt{.venv/bin/python relatorio/testes\_professor.py}
        (resultados completos em \texttt{relatorio/dados/resultados.json});
  \item Inputs patológicos gerados: \texttt{relatorio/dados/input\_*.yaml/.json}.
\end{itemize}

\end{document}
